\section{Proofs for successive cuts}\label{cv:appendix}

\subsection{Ball inclusions for the bilinear rules}\label{cv:ball-proofs}

We give the complete set formulas used in \cref{cv:rule-steps}. For the
side $S_+=\{w\le xy\}$, define
\[
 \widehat x_+(s)=\left(\frac{x-y}{2},\frac{w+1}{2}\right),\qquad
 \widehat y_+(s)=\left(\frac{x+y}{2},\frac{w-1}{2}\right).
\]
For $S_-=\{w\ge xy\}$, define
\[
 \widehat x_-(s)=\left(\frac{x+y}{2},\frac{1-w}{2}\right),\qquad
 \widehat y_-(s)=\left(\frac{x-y}{2},\frac{-1-w}{2}\right).
\]
Then $\norm{\widehat x_\sigma(s)}^2-
\norm{\widehat y_\sigma(s)}^2$ equals the signed violation
$w-xy$ or $xy-w$. Fix the violated side, suppress its subscript, and put
\[
 a=\norm{\widehat x(\sbar)},\quad
 b=\norm{\widehat y(\sbar)},\quad
 \lambda^{\mathrm S}=\frac{\widehat x(\sbar)}{a},\quad
 \ell=\lambda^{\mathrm S}_2.
\]
Here $a>b\ge0$, since $a^2-b^2=q>0$.
The unit vector $\lambda^{\mathrm S}$ is distinct from the ray coordinates.
The analytical default Case-4 set is
\begin{equation}\label{cv:scip-definition}
 C_{\mathrm S}=\{s:\phi_\ell(\widehat y(s))
                            \le(\lambda^{\mathrm S})^T\widehat x(s)\},
 \qquad
 \phi_\ell(y)=
 \begin{cases}
 \norm{y},&y_2\le\ell\norm{y},\\
 \sqrt{1-\ell^2}\,\abs{y_1}+\ell y_2,&y_2>\ell\norm{y}.
 \end{cases}
\end{equation}
Indeed, $\phi_\ell$ is the support function of the compact convex set
$\{z:\norm{z}\le1,\ z_2\le\ell\}$: either the norm maximizer is
permitted, or the maximizer lies on $z_2=\ell$. Thus $\phi_\ell$ is
convex and $\phi_\ell(y)\le\norm{y}$. The set $C_{\mathrm S}$ is closed
and convex and contains
\[
 C_0=\{s:\norm{\widehat y(s)}
                        \le(\lambda^{\mathrm S})^T\widehat x(s)\}.
\]

For completeness, these formulas also show directly that
$C_{\mathrm S}$ is free of the violated side. Consider side $+$ first.
The default vector cannot have $\ell=-1$: that would require
$\bar x=\bar y$ and $\bar w<-1$, contradicting
$\bar w-\bar x\bar y>0$. For $-1<\ell<1$, minimizing
\[
 \norm{(y_1,y_2-t)}+\ell t\quad\hbox{over }t\ge0
\]
gives exactly $\phi_\ell(y)$, and the minimum is attained. This is the
condition for $s-2te_w$ to belong to $C_0$, so
$C_{\mathrm S}=C_0+\Rplus e_w$. For $\ell=1$ the minimum occurs at zero,
and $C_{\mathrm S}=C_0$, which is already upward closed. Every point of
$C_0$ satisfies
$\norm{\widehat y}\le(\lambda^{\mathrm S})^T\widehat x
\le\norm{\widehat x}$, and hence $w-xy\ge0$.
Moving upward preserves this inequality. Since its zero set has empty
interior and $\partial(w-xy)/\partial w=1$, an interior point of
$C_{\mathrm S}$ has $w-xy>0$. For side $-$, the substitution
$(x,y,w)\mapsto(x,-y,-w)$ reduces the assertion to side $+$.

Both affine coordinate maps have linear-part norm at most $1/\sqrt2$;
for example,
\[
 \norm{\Delta\widehat x_+}^2
 =\frac{(\Delta x-\Delta y)^2+\Delta w^2}{4}
 \le\frac{\norm{\Delta s}^2}{2}.
\]
For $\norm{s-\sbar}\le t$, it follows that
\[
 (\lambda^{\mathrm S})^T\widehat x(s)\ge a-t/\sqrt2,
 \qquad \norm{\widehat y(s)}\le b+t/\sqrt2.
\]
Consequently $C_0$, and hence $C_{\mathrm S}$, contains the closed ball
with radius
\begin{equation}\label{cv:scip-ball}
 \rho_{\mathrm S}=
 \frac{a-b}{\sqrt2}=\frac{q}{\sqrt2(a+b)}.
\end{equation}
On $[-B,B]^3$, each of $a$ and $b$ is at most
$L_B=(B^2+(B+1)^2/4)^{1/2}$. This proves
\cref{cv:rule-steps}(a), including the interior condition at the vertex.

For the orbit inclusion, set $M=M_\sigma$ as in
\cref{cv:rule-steps} and choose $F^T=M(\sbar)^{-1}$. Its determinant is
$1/q>0$. The linear part satisfies
\[
 M_{0,+}(\Delta s)=
   \begin{pmatrix}\Delta w&\Delta x\\\Delta y&0\end{pmatrix},\qquad
 M_{0,-}(\Delta s)=
   \begin{pmatrix}\Delta x&\Delta w\\0&\Delta y\end{pmatrix},\qquad
 \norm{M_{0,\sigma}(\Delta s)}_F=\norm{\Delta s}.
\]
If $\norm{\Delta s}\le\sigma_{\min}(M(\sbar))$, then
\[
 \norm{M(\sbar)^{-1}M_{0,\sigma}(\Delta s)}_2\le1,
\]
and therefore
\[
 \sym(F^TM(\sbar+\Delta s))
   =I+\sym(M(\sbar)^{-1}M_{0,\sigma}(\Delta s))\succeq0.
\]
Thus the orbit set contains the ball of radius
$\sigma_{\min}(M(\sbar))$. The determinant and singular-value identities
give
\begin{equation}\label{cv:orbit-ball}
 \sigma_{\min}(M(\sbar))
 =\frac{q}{\sigma_{\max}(M(\sbar))}
 \ge\frac{q}{\norm{M(\sbar)}_F}
 \ge\frac{q}{\sqrt{3B^2+1}}.
\end{equation}
There is no need to assume this orbit set is maximal. To check that it is
free of the violated side, an interior point must have
$\sym(F^TM(s))\succ0$: a zero quadratic form at a singular positive
semidefinite matrix would give a supporting affine inequality that is
strict at $\sbar$. For a real $2\times2$ matrix $A$, writing its symmetric
and skew parts shows
$\det A=\det(\sym A)+h^2$ for some real $h$. Hence a positive definite
symmetric part implies $\det(F^TM(s))>0$, and so $\det M(s)>0$.
This proves \cref{cv:rule-steps}(b).

Finally the reference set in (b) has
$\alpha_j\norm{r_j^r}\ge\rho_A=q/T_B$. For geometric weights its
objective is at least $\rho_A$, so a $\theta$-approximate maximizer has
$\delta\ge\theta\rho_A$. For perturbed weights its objective is at least
$\tau\rho_A$, and each finite returned step satisfies
\[
 \alpha_j\norm{r_j^r}
   =\frac{u_j\alpha_j}{\widehat\omega_j+\tau}
   \ge\frac{\theta\tau\rho_A}{1+\tau}.
\]
This proves \cref{cv:rule-steps}(c).

\subsection{Proof of the suboptimal-limit construction}\label{cv:stalling-proof}

We prove \cref{cv:stalling} without computational verification. Since
$xy=w\ge1$, AM--GM gives
\[
 \varepsilon x+y+w
 \ge2\sqrt{\varepsilon xy}+w
 =2\sqrt{\varepsilon w}+w\ge1+2\sqrt\varepsilon.
\]
The allowed point $(1/\sqrt\varepsilon,\sqrt\varepsilon,1)$ attains this
bound. With the one-sided constraint $w\le xy$, the same lower bound
holds because $xy\ge w\ge1$.

Each $C_a$ is a closed, full-dimensional convex set, and
\[
 \frac{(ax+y/a)^2}{4}-xy=\frac{(ax-y/a)^2}{4}\ge0
\]
proves that its interior avoids $\{w\le xy\}$. The invertible map
$T_a(x,y,w)=(x/a,ay,w)$ preserves $xy$ and maps $C_1$ onto $C_a$.
To prove maximality it therefore suffices to consider $a=1$. In coordinates
$u=(x+y)/2$, $v=(x-y)/2$, we have
\[
 C_1=\{w\ge u^2\},\qquad S=\{w\le u^2-v^2\}.
\]
Let $p\notin C_1$, choose $0<\tau<1$, and take
\[
 u_c=u_p,\qquad v_c=-\frac{\tau v_p}{1-\tau},\qquad
 w_c=u_p^2+\eta
\]
with $\eta>0$ sufficiently small. Then $c\in\intt C_1$, and
$z=(1-\tau)c+\tau p$ has $u_z=u_p$, $v_z=0$, and
\[
 u_z^2-v_z^2-w_z
   =\tau(u_p^2-w_p)-(1-\tau)\eta>0.
\]
It is in $\intt S$. A ball around $c$ inside $C_1$ scales to a ball
around $z$ inside $\conv(C_1\cup\{p\})$, so $z$ is also an interior
point of that convex hull. Every strict convex enlargement of $C_1$
therefore fails to be $S$-free.

At $(0,0,1)$, the default vector of \cref{cv:scip-definition} is
$\lambda^{\mathrm S}=(0,1)$. Its set is
\[
 \sqrt{(x+y)^2+(w-1)^2}\le w+1,
\]
which is equivalent to $w\ge(x+y)^2/4$. Thus it is $C_1$, and the maps
$T_a$ give the asserted images.

At round zero the box has the unique optimum $(0,0,1)$ and basis rays
$e_x,e_y,e_w$. With $a=2/\xi_1$, the steps of $C_a$ are
$\xi_1,4/\xi_1,\infty$, so the first cut is
$x/\xi_1+\xi_1y/4\ge1$. Suppose inductively that $P_r$ is the box
intersected with
\begin{equation}\label{cv:successive-halfspaces}
 \frac{x}{\xi_k}+\frac{\xi_ky}{4}\ge1,\qquad 1\le k\le r.
\end{equation}
For $r\ge1$, put $\xi=\xi_r$. The point $(\xi,0,1)$ satisfies every
row, with every earlier cut and every box row other than $y\ge0$ and
$w\ge1$ slack. The three tight rows define the basis rays
\begin{equation}\label{cv:stalling-rays}
 e_x,\qquad p_\xi=(-\xi^2/4,1,0),\qquad e_w.
\end{equation}
The objective's directional costs are
\[
 \varepsilon,\qquad1-\varepsilon\xi^2/4,\qquad1,
\]
all strictly positive because $\xi<2/\sqrt\varepsilon$. The polytope
lies in this cone, so $(\xi,0,1)$ is its unique optimum.

Write $\xi'=\xi_{r+1}$ and $a=2/\xi'$. The selected set contains the
vertex in its interior because $(a\xi)^2/4=(\xi/\xi')^2<1$.
The step on $e_x$ is $\xi'-\xi$, and that on $e_w$ is infinite.
Along $p_\xi$, the expression inside the square is
\[
 ax(t)+y(t)/a
    =\frac{2\xi}{\xi'}+
                       \frac{\xi'^2-\xi^2}{2\xi'}t.
\]
It increases from a value in $[0,2)$, and reaches two first at
$t=4/(\xi+\xi')$. In cone coordinates
\[
 \lambda_x=x-\xi+\xi^2y/4,\qquad
 \lambda_p=y,\qquad\lambda_w=w-1,
\]
the new cut is
\[
 \frac{\lambda_x}{\xi'-\xi}
       +\frac{\xi+\xi'}{4}\lambda_p\ge1.
\]
Multiplication by $\xi'-\xi$ simplifies it to
$x/\xi'+\xi'y/4\ge1$. This also covers round zero by setting
$\xi=0$. The next point $(\xi',0,1)$ satisfies all old rows and the new
cut. Its tight basis has the same form, with $\xi'$ in place of $\xi$,
and the same strictly positive directional-cost argument proves unique
optimality. This completes the induction. Every point has $w-xy=1$, and
its value converges to $1+\varepsilon\xi_\infty<1+2\sqrt\varepsilon$.

It remains to check uniform pointedness. Set
$T=\xi_\infty^2/4$ and $h=(1,T+1,1)$. For the normalized rays in
\cref{cv:stalling-rays}, write $t=\xi^2/4\in[0,T]$. Their inner products
with $h$ are
\[
 1,\qquad\frac{T+1-t}{\sqrt{1+t^2}},\qquad1,
\]
each at least $1/\sqrt{1+T^2}$. Every convex combination of the normalized
rays consequently has norm at least
\begin{equation}\label{cv:stalling-pointedness}
 \gamma_r\ge
       \frac{1}{\sqrt{1+T^2}\sqrt{2+(T+1)^2}}>0.
\end{equation}
The argument includes the initial orthant. On the other hand the new
vertex is in the retained part of the old cone, so
\[
 d_r\le\norm{v^{r+1}-v^r}=\xi_{r+1}-\xi_r\longrightarrow0.
\]
The accumulation point $(\xi_\infty,0,1)$ is infeasible. Hence the
uniform depth condition fails while uniform pointedness holds.
